\scanpage{89}


On peut donc interpréter $E_1*E_2$ comme l'ext. de $\underline O$ par $L_1\otimes L_2$ déduite de l'ext. $E_1\overset{\leftrightarrows}{\otimes} E_2$ de $O$ par $L_1\times L_2$,
\unsure{ou ce qui est la même chose|cela est bien encore|via cet homom., et}, $P_1*P_2$ comme le torseur déduit du $(L_1\times L_2)$-torseur $P_1\times P_2$ par l'homomorphisme précédent.

Retour sur \unsure{les|le} \deleted{produit} rev. quadratique.\textsuperscript{*}
% slanted

\margin{* Dire d'abord \unsure{que, pour l'instant, sur une extension|ce que l'on entend par une extension} $0\longrightarrow L\longrightarrow E\longrightarrow\underline O\longrightarrow0$ $2$-trivialisée, avec $L$ inv., et une forme quadratique $Q$ ``admissible'' $\inserted{E\mathbin{-}L^{\otimes2}}$ (pour cette structure), -- i.e. a) $Q(u)=u^2$ si $u\in\Gamma L$ b) $T_Q=Q$, i.e. $Q$ de la forme $Q(x)=x^2-Tx\psi(x)+N\psi(x)^2$, où $N\in\operatorname{Sym}^2(E)$ convenable (\unsure{après} tel que $\forall x$, $Q(x)$ soit \unsure{d'év.} dans $L^{\otimes2}$). En termes d'un scindage de $E$, par \unsure{un|une} $x\in\Gamma P$, d'où $b=\pi_T(x)\in\Gamma L$, $Q$ doit être de la forme $Q(\underset{\substack{\rotatebox{-90}{$\in$}\\\Gamma O}}{\lambda}x+\underset{\substack{\rotatebox{-90}{$\in$}\\\Gamma L}}{u})=u^2+\lambda bu+\lambda^2c$ où $c\in\Gamma L^{\otimes2}$ -- le choix d'une $Q$ admissible revient donc à celui de $c\in\Gamma L^{\otimes2}$. \unsure{D'où} les $Q$ \unsure{admissibles} forment les sections d' un torseur sous $L^{\otimes2}$.}

Supposons \deleted{\unsure{à}} que $E_1,E_2$ soient associés à deux rev. quadratiques $X_1,X_2$ de $X$.

Cela signifie donc qu'on se donne sur $E_1,E_2$ formes quadratiques $Q_1,Q_2$

\[ \begin{aligned}Q_1:E_1&\mathbin{\unsure{-|\to}}L_1^{\otimes2}\\ Q_2:E_2&\mathbin{-}L_2^{\otimes2}\end{aligned} \]

dont les restrictions à $L_1,L_2$ sont $u\mapsto u^2$.

Ou encore que cela \unsure{r. à|R. à|$R_i$} des fonctions pol. de degré $\leq2$

\[ \begin{aligned}P_1&\longrightarrow L_1^{\otimes2},\\ P_2&\longrightarrow L_2^{\otimes2},\end{aligned} \]

dont les parties homogènes soient $u\longmapsto u^2$.

La donnée de $Q_1,Q_2$ définit des discriminants
% archive sheet number: 88
